% ===================================================================
%  HOTO Na 5 — Gida da ginin sa.
%
%  Traduction, faite en 2026, en haoussa d'aujourd'hui, sur l'ido de
%  texto/io. Regles de la colonne : voir 10-hoto-01.tex. Une seule
%  numerotation. Le tableau est un DIALOGUE d'un bout a l'autre :
%  trente-six attributions de parole, le plus grand nombre du livret.
%
%  C'EST LE TABLEAU QUI AVAIT OBLIGE kolonoj.py A DISTINGUER LE RECIT
%  DU DIALOGUE, ET LA COLONNE HAOUSSA N'EN A PAS BESOIN. Le javanais
%  a deux lexiques paralleles, ngoko et krama, et un enfant qui parle
%  a son oncle DOIT passer au krama : huit regles avaient donc du
%  quitter la cle « mot » pour la cle « narracio », qui ne s'applique
%  qu'aux fichiers narres. Le haoussa n'a pas de niveaux de langue —
%  il marque le rapport dans le VOCATIF, non dans le lexique. Ioannes
%  dit « Baffana », mon oncle, et jamais son nom ; l'oncle repond
%  « ɗan uwana ». Le mot change de personne, pas de forme. Aucune des
%  douze regles de cette colonne n'a donc a se deplacer, et le fichier
%  se controle comme les quinze autres.
%
%  LE FORGERON (125) ET LA FORGE (126) SONT DANS LE MEME ALINEA, ET
%  CE SONT LES DEUX MOTS ANNONCES AU TABLEAU 1. maƙeri l'homme,
%  maƙera le lieu : meme prefixe ma-, meme radical, et le lieu se
%  separe de l'homme sur la derniere voyelle. L'en-tete du tableau 1
%  citait la paire comme exemple ; la voila dans le texte, a douze
%  mots l'une de l'autre. Et l'alinea en ajoute deux du meme moule :
%  masaba l'enclume (127) et matsi les tenailles (128) — l'outil se
%  fabrique avec le prefixe de l'homme et du lieu.
%
%  CE TABLEAU EST LE PLUS « ma- » DU LIVRET, ET IL FALLAIT S'Y
%  ATTENDRE : IL NE PARLE QUE DE METIERS ET D'OUTILS. On y compte
%  magini le macon (108), masassaƙi le charpentier (115), maƙeri le
%  forgeron, masanin gini l'architecte (75), ma’aikacin hannu le
%  manoeuvre (146), makulli la serrure (6), matsefin lambu le rateau
%  (130), ma’aunin tsaye le fil a plomb (109), magudanar rufi la
%  gouttiere (42), matakala l'escalier (13), matakai les marches
%  (14), masaba, matsi, mariƙi. Une seule syllabe fait tout le
%  chantier.
%
%  ET TROIS EMPRUNTS DE CE TABLEAU MONTRENT LA LOI DU p AUSSI BIEN
%  QUE LE TABLEAU 1 : famfo la pompe (58), fenti la peinture (103),
%  falo le salon (b), de parlour. Le boko n'a pas de p, et il n'en
%  laisse pas passer un seul, meme dans un mot venu la semaine
%  derniere.
%
%  TROISIEME REFUS DE LA COLONNE : LA BRIQUE (61). Le haoussa a
%  « tubali », et c'est le mot le plus tentant du tableau : le tubali
%  est la brique de terre crue, moulee a la main en forme de poire,
%  et ce sont des tubali qui ont bati les murailles de Kano et les
%  mosquees a colonnes du Sahel. La planche grave des briques cuites,
%  parallelepipedes, sorties d'un four. La colonne ecrit « bulo ». Un
%  mot d'ici ne se pose pas sur une chose de la-bas — apres le bori
%  du tableau 3 et le pardessus du tableau 4, c'est le troisieme.
%
%  QUATRIEME REFUS, AU DERNIER ALINEA : LE GRENIER (34). « rumbu »
%  est le grenier haoussa, et il n'est pas une piece : c'est un grand
%  vase de terre et de paille tressee, pose dans la cour, ou le mil
%  se garde d'une recolte a l'autre. Le grenier de la planche est un
%  comble, sous les ardoises. On ecrit « ɗakin ajiyar hatsi », la
%  chambre ou l'on garde le grain.
%
%  DEUX MOTS QUI SONT DEUX CHOSES, ET LES DEUX PARAISSENT ICI.
%  « kwano » est le plat — celui de la tisane au tableau 4 — et c'est
%  aussi la TOLE ONDULEE (24), qui couvre aujourd'hui presque tous
%  les toits du Nigeria ; le zingueur (121) est « mai aikin kwano ».
%  « allo » est l'ardoise de l'ecolier au tableau 1 et l'ardoise du
%  toit (66) ici : meme pierre, meme mot, et le haoussa n'a pas
%  eprouve le besoin de les separer. Apres « kaka » du tableau 4,
%  c'est la troisieme paire d'homonymes que ce releve rencontre dans
%  cette langue.
%
%  UNE EXEMPTION, ET NON UN RESSERREMENT. La note « Balk-o » de ce
%  tableau est le SEUL endroit du livret qui cite d'autres langues mot
%  pour mot — Balken, joist, solive, travicello, balka, viga —, et
%  l'ido les met lui-meme en italique pour dire qu'elles sont
%  etrangeres. La regle des lettres absentes du boko a donc crie
%  quatre fois, sur solive, poutre et viga. Elle avait raison a la
%  lettre : ces mots ne sont pas du haoussa et ne doivent surtout pas
%  le devenir. On n'a donc RIEN resserre — une citation n'est pas une
%  faute d'orthographe, c'est un autre texte — et les trois lignes qui
%  la portent sont passees en « exemptes », comme la colonne arabe
%  egyptienne exempte la plaisanterie du vieux cordonnier. Verifie
%  apres coup : « pompa » glisse dans un alinea ordinaire est releve.
%
%  ET LE CABINET D'AISANCES (20) SE DIT « bayan gida », DERRIERE LA
%  MAISON. Le mot ne nomme pas la chose mais sa PLACE, et il la nomme
%  encore quand la chose a change de place — le plan de ce tableau en
%  met un a l'etage. C'est la meme facon de faire que le vasistas du
%  tableau 1, « ƙaramar taga », et que le maƙogwaro : le haoussa
%  nomme volontiers par ou plutot que par quoi.
% ===================================================================

%%K t05-apar-1 apar
\VUsaut{6.0mm}
\VUtitre{15.0pt}{60}{HOTO Na 5}
\VUsaut{1.4mm}
\VUtitre{11.4pt}{0}{Gida da ginin sa.}
\VUsaut{2.2mm}

%%K t05-tit-1 sub
\VUpk{10.2pt}{40}{Kyakkyawar haɗuwa.}
\VUsaut{3.0mm}

%%K t05-01-1 p
\textsc{Ioannes}. --- Baffana, ina son sanin yadda ake gina
\VUgras{gida.}

%%K t05-01-2 p
\textsc{Baffa} (80). --- Wannan abu ne mai sauƙi ƙwarai, ɗan uwana, zo
tare da ni, zan nuna maka wanda malam Bernardus (79) yake sa a gina
kuma wanda zan zauna a ciki.

%%K t05-01-3 p
\textsc{Ioannes}. --- Wa ya zana \VUgras{shirin gini}
\textsuperscript{(76)} na wannan gida ?

%%K t05-01-4 p
\textsc{Baffa}. --- \VUgras{Masanin gini} \textsuperscript{(75)} ; ya
gabatar da zane bayyananne ƙwarai wanda aka amince da shi, sai
\VUgras{ɗan kwangila} \textsuperscript{(77)} ya fara aikin.

%%K t05-01-5 p
\textsc{Ioannes}. --- Wane ne kuwa ɗan kwangilar ?

%%K t05-01-6 p
\textsc{Baffa}. --- Shi ne malam Blanko, mahaifin abokinka Iakobus.
Yana jagorantar ma’aikata tare da malam Rufo, masanin gini.

%%K t05-01-7 p
Ga shi ! Malam Blanko yana zuwa a daidai lokaci da \VUgras{rularsa}
\textsuperscript{(78)}, malam Rufo kuma da shirin ginin sa.

%%K t05-01-8 p
Barka da rana, malamai. Ina lafiya ?

%%K t05-01-9 p
\textsc{Malam Rufo}. --- Lafiya lau, na gode. Barka da rana, Ioannes,
kai, yaron nan ya yi girma !

%%K t05-01-10 p
\textsc{Baffa}. --- I, hakika, ƙaramin mutum ne. Yana da natsuwa
ƙwarai, dole a sanar da shi kome da ya gani. Yau, yana son ya san
yadda ake gina gida.

%%K t05-tit-2 sub
\VUpk{10.2pt}{40}{Ma’aikata.}
\VUsaut{3.0mm}

%%K t05-01-11 p
\textsc{Malam Blanko}. --- To, zo tare da ni, ƙaramin aboki, zan
bayyana maka nan take, domin ina son yara masu neman ilimi ƙwarai.

%%K t05-01-12 p
Ka ga ma’aikacin nan da yake riƙe da \VUgras{shebur}
\textsuperscript{(82)} a hannu : shi ne \VUgras{mai haƙa ƙasa}
\textsuperscript{(81)}. Yana haƙa \VUgras{rami} \textsuperscript{(46)}
domin a sa bututun ruwa. Ya kuma haƙa ƙasar a inda gidan yake, domin
magini ya kafa \VUgras{bangaye} huɗu. Ɓangaren waɗannan bangaye da
yake cikin ƙasa ana kiransa \VUgras{harsashi} \textsuperscript{(2)}.
Wurin da yake tsakanin harsashin yana yin \VUgras{ɗakin ƙasa}
\textsuperscript{(3)}. Wannan ɗakin ƙasa yana da ƙaramar taga da ake
kira \VUgras{tagar ɗakin ƙasa} \textsuperscript{(5)}, da \VUgras{ƙofa}
\textsuperscript{(4)} da matakala.

%%K t05-01-13 p
\VUgras{Magini} \textsuperscript{(108)} ya ci gaba da gina bangayen
yana bar buɗaɗɗun wurare domin \VUgras{ƙofofi} \textsuperscript{(8)} da
\VUgras{tagogi} \textsuperscript{(17)}.

%%K t05-01-14 p
\textsc{Ioannes}. --- Amma ban ga maginin nan ba !

%%K t05-01-15 p
\textsc{Malam Blanko}. --- Yanzu ya gama aiki a cikin gidan. Yana kusa
da \VUgras{gidan dawaki} \textsuperscript{(54)}, a kan \VUgras{matakalar
aikinsa} \textsuperscript{(110)}, yana gina bangon \VUgras{wurin wanki}
\textsuperscript{(56)}.

%%K t05-01-16 p
Yana da cokalin magini, da kwanon laka da \VUgras{ma’aunin tsaye}
\textsuperscript{(109)}. Amma ba za ka tuna sunayen nan ba.

%%K t05-01-17 p
\textsc{Ioannes}. --- Lalle zan tuna, domin nan take zan roƙi baffana
ƙaramin cokalin magini da ƙaramin kwanon laka, ni kuma da kaina zan yi
ma’aunin tsaye da igiya.

%%K t05-tit-3 sub
\VUsaut{3.0mm}
\VUpk{10.2pt}{40}{Kayan gini.}
\VUsaut{3.0mm}

%%K t05-01-18 p
\textsc{Baffa}. --- Kai ! madalla. Amma gaya mini, Ioannes, me ake
buƙata domin a gina gida ?

%%K t05-01-19 p
\textsc{Ioannes}. --- Ana buƙatar \VUgras{duwatsun gini}
\textsuperscript{(59)} da \VUgras{laka} \textsuperscript{(65)}.

%%K t05-01-20 p
\textsc{Baffa}. --- Daidai, ka ga mutumin nan mai babbar hula, a kan
\VUgras{karusarsa} \textsuperscript{(136)}. Shi ne \VUgras{mai tuƙa
karusa} \textsuperscript{(135)}. Kowace rana yana kawo nan \VUgras{manyan
duwatsu} \textsuperscript{(60)}. Yana saukar da kayansa a farfajiya,
\VUgras{masu sassaƙa dutse} \textsuperscript{(83)} kuma suna sassaƙa
duwatsun suna kuma yanka su da \VUgras{zarton dutse}
\textsuperscript{(85)}. \VUgras{Ma’aikacin hannu} \textsuperscript{(146)}
yana kai su ga magini wanda yake ajiye su.

%%K t05-01-21 p
A halin yanzu, \VUgras{mai gauraya} \textsuperscript{(87)} yana shirya
laka. Yana buƙatar \VUgras{yashi} \textsuperscript{(62)}, da \VUgras{farar
ƙasa} \textsuperscript{(63)} da \VUgras{ruwa} \textsuperscript{(64)}.
Farar ƙasar tana kusa da shi cikin \VUgras{buhu} \textsuperscript{(71)},
\VUgras{mai taimakon magini} \textsuperscript{(111)} kuma yana kawo masa
yashi a kan \VUgras{amalanke} \textsuperscript{(112)}, \VUgras{ɗan koyo}
\textsuperscript{(113)} kuwa yana kan famfo (58). Yana cika
\VUgras{guga} \textsuperscript{(114)}. Ba da daɗewa ba lakar za ta yi.
\VUgras{Mai ɗaukar laka} \textsuperscript{(107)} yana ɗaukarta yana hawa
da ita, ta tsani, zuwa kan matakalar aiki, inda magini yake aiki yana
gama wannan gefen gidan.

%%K t05-01-22 p
Yanzu kuma, yaya ake kiran ma’aikacin da yake aikin \VUgras{rufi}
\textsuperscript{(31)} ?

%%K t05-01-23 p
\textsc{Ioannes}. --- Shi ne \VUgras{mai rufin gida}
\textsuperscript{(118)}.

%%K t05-tit-4 sub
\VUsaut{3.0mm}
\VUpk{10.2pt}{40}{Rufin gida.}
\VUsaut{3.0mm}

%%K t05-01-24 p
\textsc{Malam Blanko}. --- I, amma da farko sai \VUgras{masassaƙi}
\textsuperscript{(115)} ya zo.

%%K t05-01-25 p
Masassaƙi yana aikin \VUgras{katakon rufi} \textsuperscript{(35)}. Yana
haɗa \VUgras{manyan katakai} \textsuperscript{(37)} da \VUgras{ƙusoshi}
\textsuperscript{(117)}. Ma’aikatan nan, can sama, masu fadaɗɗun
gajerun wandunan kaɗi, masassaƙa ne. Ɗaya yana riƙe da ƙaramin
katako, ɗayan kuma yana yanka ƙusa da \VUgras{gatarinsa}
\textsuperscript{(116)}. Wanda yake kusa da \VUgras{bututun hayaƙi}
\textsuperscript{(41)} shi ne mai rufin gida. Yana ƙusa sandunan rufi a
kan (*) katakan bene, da \VUgras{allunan dutse} \textsuperscript{(66)} a
kan sandunan. Wani lokaci yakan sa tayal.

%%K t05-noto-1 noto
\VUnotes{143.2mm}{%
(*) \VUgras{Balk-o} = D. \textit{Balken}, E. \textit{joist}, F.
\textit{solive}, I. \textit{travicello}, R. \textit{balka}, S. Port.
\textit{viga}.
Saiwar sana’a wadda ba a karɓe ta ba har yanzu, amma ya kamata a
gabatar da ita domin fassara ma’anar kalmar F. \textit{solive}, wadda
ba trabo ba ce, F. \textit{poutre}, ba kuwa trabeto ba. Sandar katako
ce mai kusurwa huɗu wadda take ɗauke da benen ƙasa da silin.%
}

%%K t05-01-26 p
Rufin gidan dawaki an rufe shi da \VUgras{tayal} \textsuperscript{(55)},
na gidaje kuwa an rufe su da \VUgras{allunan dutse}
\textsuperscript{(32)}, na baranda kuma da \VUgras{kwano}
\textsuperscript{(24)}.

%%K t05-01-27 p
\textsc{Ioannes}. --- Shin mai rufin gida shi ne yake aikin rufin
kwano ma ?

%%K t05-01-28 p
\textsc{Malam Blanko}. --- A’a, sai \VUgras{mai aikin kwano}
\textsuperscript{(121)} ko \VUgras{mai aikin dalma,} shi ne yake sa
\VUgras{tagar rufi} \textsuperscript{(38)} a kan rufin gidan. Yana sa
\VUgras{bututun ruwa} \textsuperscript{(43)} da \VUgras{magudanar rufi}
\textsuperscript{(42)} ma. Yana walƙa kwano da farar tama. Shi ne ya
kafa \VUgras{sandar tsawa} \textsuperscript{(40)} da \VUgras{alamar
iska} \textsuperscript{(39)} a kan \VUgras{gaban rufi}
\textsuperscript{(36)}.

%%K t05-01-29 p
\textsc{Ioannes}. --- Haka nan gidan ya cika kenan ?

%%K t05-tit-5 sub
\VUsaut{3.0mm}
\VUpk{10.2pt}{40}{Kayan aiki.}
\VUsaut{3.0mm}

%%K t05-01-30 p
\textsc{Malam Blanko}. --- Bai cika sarai ba. \VUgras{Mai shafe
bango} \textsuperscript{(99)} yana gina bangayen da suke raba shi zuwa
ɗakuna iri iri da \VUgras{bulo} \textsuperscript{(61)}. Yana shafa
\VUgras{farar siminti} \textsuperscript{(70)} a kan \VUgras{silin}
\textsuperscript{(26)} da bangaye. Yanzu yana aikin silin na
\VUgras{baranda} \textsuperscript{(33)}. Yana da, kamar magini,
\VUgras{cokalin magini} \textsuperscript{(100)} da \VUgras{kwanon laka}
\textsuperscript{(101)}.

%%K t05-01-31 p
Sai kafinta ya yi aikin ƙofofi, da tagogi, da \VUgras{katakon ƙasa}
\textsuperscript{(16)} da \VUgras{murfin taga} \textsuperscript{(19)}.

%%K t05-01-32 p
Yanzu yana aiki a farfajiya a kan \VUgras{teburin aikinsa}
\textsuperscript{(90)}. Yana da \VUgras{zarto} \textsuperscript{(94)}
domin ya yanka katako (69), da \VUgras{abin santsin katako}
\textsuperscript{(91)}, da \VUgras{mariƙin katako}
\textsuperscript{(92)}, da \VUgras{ƙarfen sassaƙa}
\textsuperscript{(94 bis)}, da \VUgras{kwamfas} \textsuperscript{(95)},
da \VUgras{gudumar katako} \textsuperscript{(95 bis)}.

%%K t05-01-33 p
\textsc{Ioannes}. --- Kai ! aikin kafinta ya fi na magini daɗi.
Baffana, za ka saya mini teburin aiki, da zarto da abin santsin
katako, domin ba da daɗewa ba zan yi wa Medor ɗaki.

%%K t05-01-34 p
\textsc{Baffa}. --- I, ɗan ƙarami.

%%K t05-01-35 p
\textsc{Ioannes}. --- Kai, ina cikin farin ciki ! Wane ne kuma wanda
yake tsaye kusa da ƙofar shiga ?

%%K t05-tit-6 sub
\VUsaut{3.0mm}
\VUpk{10.2pt}{40}{Shafa fenti.}
\VUsaut{3.0mm}

%%K t05-01-36 p
\textsc{Malam Blanko}. --- Shi ne \VUgras{mai shafa fenti}
\textsuperscript{(102)}. Yana tsaye a kan tsaninsa da kwanon
\VUgras{fenti} \textsuperscript{(103)} da \VUgras{buroshinsa}
\textsuperscript{(104)}. Yana shafa katakon ƙofar shiga domin ya daɗe.

%%K t05-01-37 p
Shi ne kuma yake liƙa takardun bango a cikin ɗakuna.

%%K t05-01-38 p
\VUgras{Mai kayan ado na gida} \textsuperscript{(107)} wanda yake kan
\VUgras{tsani} \textsuperscript{(106)}, a kan \VUgras{barandar sama}
\textsuperscript{(28)}, yana sa \VUgras{labule}
\textsuperscript{(29)}.

%%K t05-01-39 p
Ma’aikacin da yake tsaye kusa da babbar taga shi ne \VUgras{mai sa
gilashi} \textsuperscript{(96)}. Yana kafa \VUgras{gilashin taga}
\textsuperscript{(18)} da \VUgras{manne} \textsuperscript{(73)}. Ɗayan
ma’aikacin, wanda yake durƙushe kusa da ƙofar ɗakin ƙasa, shi ne
\VUgras{mai yin makulli} \textsuperscript{(97)}. Yana sa
\VUgras{makulli} \textsuperscript{(6)}. \VUgras{Abin nika ƙarfensa}
\textsuperscript{(98)} yana kusa da shi.

%%K t05-01-40 p
\textsc{Ioannes}. --- Shin mai yin makulli ne shi ma, ma’aikacin nan
da yake bugu da \VUgras{gudumarsa} \textsuperscript{(84)} a kan ƙarfe ?

%%K t05-01-41 p
\textsc{Malam Blanko}. --- Daidai da gaskiya, maƙeri ne (125). Yana
ƙera \VUgras{kayan ƙarfe} \textsuperscript{(67)} domin \VUgras{mai wutar
lantarki} \textsuperscript{(124)} wanda yake kan rufin \VUgras{gidan
karusa} \textsuperscript{(51)}. Yana da \VUgras{maƙera}
\textsuperscript{(126)}, da \VUgras{masaba} \textsuperscript{(127)} da
\VUgras{matsi} \textsuperscript{(128)}.

%%K t05-01-42 p
Mai wutar lantarki yana kafa \VUgras{wayar} \VUgras{jan ƙarfe}
\textsuperscript{(44)} ta tarho.

%%K t05-tit-7 sub
\VUpk{10.2pt}{40}{Aikin lambu.}
\VUsaut{3.0mm}

%%K t05-01-43 p
\textsc{Baffa}. --- A ƙarshen \VUgras{farfajiya} \textsuperscript{(45)},
kusa da \VUgras{shingen ƙarfe} \textsuperscript{(47)} da \VUgras{babbar
ƙofa} \textsuperscript{(48)}, ka ga \VUgras{mai lambu}
\textsuperscript{(129)}. Yana shirya \VUgras{ƙaramin lambu}
\textsuperscript{(49)}. Ya huda ƙasar da \VUgras{fartanyarsa}
\textsuperscript{(131)}, yana shuka iri yana kuma share da
\VUgras{matsefin lambu} \textsuperscript{(130)}. Sa’an nan, da
\VUgras{abin shayarwarsa} \textsuperscript{(132)}, zai shayar da
\VUgras{gadajen shuka} \textsuperscript{(150)}, tsire-tsire kuwa za su
yi girma.

%%K t05-01-44 p
\VUgras{Mai kula da doki} \textsuperscript{(133)} wanda ya gama kula da
dokin ya kuma ba shi abinci, yana wanke \VUgras{keken doki}
\textsuperscript{(52)} da soso, da \VUgras{famfon hannu}
\textsuperscript{(134)} da guga na ruwa. Sa’an nan zai mayar da shi
cikin gidan karusa ya rufe ƙofa da \VUgras{sakata}
\textsuperscript{(53)} mai kauri.

%%K t05-01-45 p
Ga shi ka gamsu, ina tsammani, ka sami isassun bayanai.

%%K t05-01-46 p
\textsc{Ioannes}. --- I, baffana, amma ina son ganin cikin gidan.

%%K t05-01-47 p
\textsc{Baffa}. --- Wannan zai faru gobe. Malam Rufo zai bayyana maka
shirin ginin ya kuma nuna maka sunan kowane ɗaki.

%%K t05-tit-8 sub
\VUsaut{3.0mm}
\VUpk{10.2pt}{40}{Tsarin ciki na gidan.}
\VUsaut{2.4mm}

%%K t05-apar-2 apar
{\centering\textit{(Duba shirin ginin.)}\par}
\VUsaut{2.4mm}

%%K t05-01-48 p
\textsc{Masanin gini}. --- Zo, Ioannes, zan sa ka bincika sassan gidan
iri iri nan take.

%%K t05-01-49 p
Mu shiga \VUgras{benen ƙasa} \textsuperscript{(7)}. Ga maɓallin
\VUgras{ƙararrawa} \textsuperscript{(11)}. Muna hawan \VUgras{matakai}
\textsuperscript{(14)} uku na \VUgras{matakalar shiga}
\textsuperscript{(9)}, muna ƙetare \VUgras{bakin} \textsuperscript{(10)}
\VUgras{ƙofar shiga} \textsuperscript{(8)}, ga mu kuwa a gindin
\VUgras{matakala} \textsuperscript{(13)}.

%%K t05-01-50 p
\textsc{Ioannes}. --- Wannan ita ce hanyar wucewa.

%%K t05-01-51 p
\textsc{Masanin gini}. --- A’a, wannan shi ne \VUgras{ɗakin shiga}
\textsuperscript{(12)}. \VUgras{Hanyar wucewa} \textsuperscript{(a)}
tana can ƙarshe. A dama, ga \VUgras{falo} \textsuperscript{(b)} da
\VUgras{ɗakin cin abinci} \textsuperscript{(c)}, a gaban ɗakin cin
abinci kuma akwai \VUgras{ɗakin girki} \textsuperscript{(d)} da
\VUgras{ɗakin ma’aikata} \textsuperscript{(e)}, sa’an nan a gaba
\VUgras{ɗakin rubutu} \textsuperscript{(f)}. \VUgras{Baranda}
\textsuperscript{(23)} tare da \VUgras{ƙofarta mai gilashi}
\textsuperscript{(25)} tana kusa da falo.

%%K t05-01-52 p
Nan take za mu hau matakala. Ka riƙe \VUgras{hannun matakala}
\textsuperscript{(15)} ka ƙidaya matakan. Su goma sha huɗu ne. Ga
\VUgras{hanyar sama} \textsuperscript{(m)}, muna kan \VUgras{bene na
farko} \textsuperscript{(27)}.

%%K t05-tit-9 sub
\VUsaut{3.0mm}
\VUpk{10.2pt}{40}{Bene na farko.}
\VUsaut{3.0mm}

%%K t05-01-53 p
A gaban matakala, akwai \VUgras{bayan gida} \textsuperscript{(20)} da
\VUgras{wurin ado} \textsuperscript{(j)}. A dama akwai kyakkyawan
\VUgras{ɗakin kwana} \textsuperscript{(g)} mai tagogi biyu a
\VUgras{gaban gidan} \textsuperscript{(1)}. A hagu, akwai \VUgras{ɗakin
wanka} \textsuperscript{(1)} da \VUgras{ɗakin baƙi}
\textsuperscript{(i)} tare da babban \VUgras{ɓoyayyen wurin gado}
\textsuperscript{(h)} : kusa da ɗakin kwana akwai \VUgras{ɗakin yara}
\textsuperscript{(k)}. Yana kusa da babban \VUgras{dandali}
\textsuperscript{(21)}. A ƙarƙashin wannan dandali akwai wani bayan
gida (20), a bango kuma, ga \VUgras{ƙararrawa} \textsuperscript{(22)}
wadda take kaɗawa domin abincin dare.

%%K t05-01-54 p
\textsc{Ioannes}. --- Mu hau \VUgras{bene na biyu.}

%%K t05-01-55 p
\textsc{Masanin gini}. --- Babu bene na biyu. A ƙarƙashin rufin akwai
\VUgras{ɗakunan rufi} \textsuperscript{(33)} kaɗai da ɗakin ajiyar
hatsi (34). Mun gama bincike, za mu iya sauka.

%%K t05-01-56 p
\textsc{Ioannes}. --- Na gode, malam, wannan yana da ban sha’awa
ƙwarai. Nan take zan ba mahaifina labarin kome da na gani.

\VUsaut{4.0mm}
\VUfilet{22mm}
